\section{Reproducibility and artifact provenance}\label{app:reproducibility}
\subsection{Delivered artifacts and numerical scope}
The accompanying source package contains the complete LaTeX manuscript, local bibliography, generated tables and figures, every individual benchmark outcome in compact JSON/CSV, all 51 generated parameter sets, and scripts that regenerate the displayed evidence. The separate frozen research supplement contains the solver and model sources, pinned environment declarations, saved schedules, raw witnesses and solver logs, conic references, analysis, and development provenance. It preserves third-party model licenses and source attributions. No solver binaries, environments or downloaded literature PDFs are redistributed. The manuscript and supplied figures build without access to the original repository, its literature directory, a network service or a solver license. The package has no claimed public archival identifier; it is supplied directly with the manuscript.

The file \texttt{README.md} gives executable instructions. The compact data description is \path{data/README.md}, its immutable-input hashes are in \path{data/provenance.json}, and its displayed entries are in \path{data/table_values.json}. The full supplement contains 9076 payload files and its embedded manifest (9077 archive members). Earlier pilots, failed exports, preliminary cone passes and superseded analyses remain as provenance; the final study uses \texttt{analysis\_v1} and the distinct eight benchmark batches identified in \cref{sec:design}. Repetitions and outcome-triggered follow-ups never replace primary outcomes.

Table regeneration redoes aggregation from the compact records, not primal validation. The separate independent audit re-evaluates witnesses in freshly constructed original GDPs, checks schedules and source metadata, and independently recomputes summary arithmetic. As stated in \cref{sec:acceptance}, it reuses the original-model checker. Acceptance uses \eqref{eq:accept-gap} and the criteria of \cref{sec:acceptance}, and cone-reference comparison uses \eqref{eq:comparison-tolerance}; these are numerical checks, not exact optimality, dual-feasibility or infeasibility certificates.

\subsection{Immutable identifiers and environments}
The following SHA-256 identifiers fix the research supplement, original executable-source freeze, and declared primary and supplementary schedules, respectively:
\begin{center}\scriptsize\ttfamily
f0399194ca1c9c57421965e236302c62e10eb72846ab807f936f18d9927d4b26\\
6c9eaf9898f04d08a12381879ab9bae79f83b800f0db9e4cc696ccbff6be84ad\\
1fd3118da6e10b5a47377ea00e3c8237262558b09764ba07ed5a66dc088090ca\\
6b45611e18668951d9794af8d7a80a8e9565a23323aee6d92cf1ad06cb4252a5
\end{center}
The primary and supplementary plan filenames are \path{study_plan_v1.json} and \path{supplementary_plan_v1.json}. Source-file and raw-record hashes, including both follow-up inputs and the post-freeze reference files, are recorded in the local provenance and full publication manifests. The publication archive is 38\,086\,636 bytes. It remains byte-identical to the completed research record; manuscript development did not alter the solver, instances or benchmark outcomes.

The benchmark software, processor, concurrency and effective solver options are specified in \cref{sec:settings}. The archived lockfiles reproduce the Python dependencies with \texttt{uv sync --frozen}; native GAMS, Gurobi and Ipopt installations and applicable licenses must be supplied separately for optimization reruns. The continuous CVXPY/Clarabel reference has its own lockfile under \texttt{lbesh\_research/conic\_reference\_env}. Historical environment declarations and native logs distinguish Gurobi 13.0.3 from GAMS/Gurobi 13.0.2. Hardware evidence comes from the saved logs, not an assumed identity with the machine used to typeset the paper. The writing environment used Python 3.13.11 and Matplotlib 3.11.2 for figure regeneration; the supplied generated files suffice for the TeX build.

\subsection{Reproduction commands and verification boundary}
From the extracted paper directory, the minimal build and local evidence checks are
\begin{verbatim}
latexmk -pdf -interaction=nonstopmode -halt-on-error main.tex
python scripts/regenerate.py
python scripts/check_evidence.py
\end{verbatim}
The generator requires Matplotlib for its two PDF figures; \texttt{--no-figures} regenerates tables and CSV using only the standard library. It verifies immutable input hashes and all 1464 unique batch--instance--method records before writing outputs. The evidence check also independently reconstructs all 51 parameter streams. The supplied source manifest and delivery checksums identify the exact manuscript version; \texttt{scripts/package.py} refreshes them after a correction.

The following verifies every research archive member against the manifest and extracts only regular, uniquely named files with safe relative paths into a new or empty directory:
\begin{verbatim}
python scripts/verify_supplement.py \
  supplement/publication_bundle_v1.tar.gz --extract research
python evidence/check_stage02.py --research-root research
\end{verbatim}
The second command checks the explicit analytic examples and selected frozen-source settings. Its research-root argument is mandatory, so relocation never silently bypasses source validation. From \texttt{research/code/minlp\_solver\_lab}, with the pinned Python environment active, the fresh full-study audit is
\begin{small}\begin{verbatim}
export PYTHONPATH="$PWD/instances/gdplib_src"
python lbesh_results_independent_audit.py \
  --fresh-validation \
  --compare-analysis results/lbesh_development/analysis_v1/analysis.json \
  --sensitivity results/lbesh_development/gurobi_trig_sensitivity_v1.jsonl \
  --legacy-initialization \
    results/lbesh_development/legacy_initialization_v1.jsonl \
  --out /tmp/lbesh-new-audit.json
\end{verbatim}\end{small}
The audit refuses an existing output file; choose a new path on each run. The explicit import path selects the extracted GDPlib sources even if a reused Python environment has an editable installation elsewhere. Set \texttt{OPENBLAS\_NUM\_THREADS}, \texttt{OMP\_NUM\_THREADS}, \texttt{MKL\_NUM\_THREADS} and \texttt{NUMEXPR\_NUM\_THREADS} to one. Optimization schedules and the separately pinned conic environment are documented in the archived \texttt{LBESH\_RESEARCH.md}; their declared commands and job orders are preserved in the plans. Executing those schedules is distinct from auditing their saved results.

As a portability check, both archives were relocated to a temporary directory, where a clean TeX build, regeneration of all 17 tables and two figures, the compact evidence check, the source/example check against the extracted research tree, and a fresh full-study audit all completed. The audit reconciled all 1464 benchmark outcomes and 6122 exported analysis fields, revalidated all 174 feasible enumeration witnesses, and found no cross-witness bound/status contradiction. Regenerated outputs matched the delivered files byte for byte. This check reused the existing pinned Python installation; it was neither a clean environment installation nor a rerun of optimization timings.
